\documentclass[../../../../main.tex]{subfiles}
\begin{document}

$(x,y)$ 平面において半円：$x^2+y^2=1$，$y \geqq 0$ の内側が鏡になっているとする．
図のように，定点 $(1,0)$ より $x$ 軸となす角 $\theta$ で光線が発射され，2回半円に反射したのち，$x$ 軸上の点 $P$ を通過したとする．

\begin{figure}[htbp]
  \centering
  \begin{tikzpicture}[scale=2, >=stealth]
    % 座標軸
    \draw[->] (-1.4, 0) -- (1.5, 0) node[right] {$x$};
    \draw[->] (0, -0.6) -- (0, 1.3) node[above] {$y$};
    \node[below left] at (0,0) {$O$};

    % 半円
    \draw[thick] (1,0) arc (0:180:1);

    % 点 (1,0)
    \coordinate (A) at (1,0);
    \node[below right] at (A) {$(1,0)$};

    % 1回目の反射点 Q1
    \coordinate (Q1) at (60:1);
    \draw[thick] (A) -- (Q1);

    % 角度 theta
    \draw (0.8, 0) arc (180:120:0.2);
    \node at (0.7, 0.18) {$\theta$};

    % 2回目の反射点 Q2
    \coordinate (Q2) at (140:1);
    \draw[thick] (Q1) -- (Q2);

    % 点 P
    \coordinate (P) at (-0.6, 0);
    \node[above left] at (P) {$P$};

    % 反射光線の延長
    \coordinate (END) at ($ (Q2)!1.4!(P) $);
    \draw[thick, ->] (Q2) -- (END);
  \end{tikzpicture}
\end{figure}

\begin{enumerate}
  \item このような状況が起こるための $\theta$ の範囲を求めよ．
  \item $P$ の座標を $\theta$ を用いて表せ．
  \item $\theta$ が (1) の範囲を動くときの $P$ の動く範囲を求めよ．
\end{enumerate}

\end{document}
